\documentclass[10pt,a4paper]{amsart}
\usepackage[margin=19mm]{geometry}
\usepackage{amsmath,amssymb,mathtools}
\usepackage[scr=rsfs,cal=euler]{mathalfa}
\usepackage{xfrac}
\usepackage[unicode,hidelinks,bookmarks=false]{hyperref}
\newcommand{\ie}{i.e.\ }
\newcommand{\cf}{cf.\ }
\newcommand{\resp}{respectively\ }
\newcommand{\Subsection}{\subsection}
\providecommand{\nobreakdash}{\nobreak\hbox{-}}
\setlength{\emergencystretch}{2em}
\multlinegap0pt
\usepackage{bm}
\usepackage{bbm}
\usepackage{dsfont}
\usepackage{xspace}
\usepackage{cancel}
\usepackage{mathtools}
\usepackage{cleveref}
\usepackage{amsthm}
\makeatletter
\@ifundefined{theorem}{\newtheorem{theorem}{Theorem}}{}
\@ifundefined{lemma}{\newtheorem{lemma}[theorem]{Lemma}}{}
\@ifundefined{proposition}{\newtheorem{proposition}[theorem]{Proposition}}{}
\makeatother
\DeclareMathOperator{\Tropw}{Trop_\mathcal{C}^+}
\DeclareMathOperator{\Trop}{Trop}
\DeclareMathOperator{\initial}{in}
\newcommand{\CC}{\mathcal{C}}
\newcommand{\RR}{\mathcal{R}}
\newcommand{\inw}{\operatorname{in}_w}
\title[Computing positive tropical varieties and lower bounds on th]{Computing positive tropical varieties and lower bounds on the number of positive roots}
\author{Kemal Rose, M\'at\'e L. Telek}
\date{Source: arXiv:2408.15719v2 (CC BY 4.0)}
\begin{document}
\maketitle

\noindent\emph{Source and licence.} Computing positive tropical varieties and lower bounds on the number of positive roots. Version arXiv:2408.15719v2 (CC BY 4.0), distributed under the Creative Commons Attribution 4.0 International licence, as recorded in the arXiv licence metadata for this exact version and shown on the abstract page https://arxiv.org/abs/2408.15719v2. Full rights notice: licenses/arxiv-2408.15719-CC-BY-4.0.txt.\par\smallskip

\noindent\textit{Editorial scope.} Counted unit: the Proposition whose source label is \texttt{prop: functoriality of monomial maps} in the article (source0.tex lines 818--833, source label \texttt{prop: functoriality of monomial maps}), delivered together with the article's own complete original proof of it (source0.tex lines 834--851, one \texttt{proof} environment of 18 lines). No mathematics is written, repaired or re-derived: statement and proof are the article's own text line for line. Required context retained uncounted: Prop::CharWeakSigned (lines 387--396); prop: positive part of toric vars (lines 765--769); lemma: induced map of initial degs (lines 784--794). Every definition, display and label of those lines is retained unchanged. Omitted from this package and not redistributed: 1--386, 397--764, 770--783, 795--817, 852--1741. Nothing inside a retained excerpt is omitted or altered: every retained body is delivered exactly as the article's lines are written, so no omission receipt of any kind accompanies this package. Citations: the retained excerpts carry no citation command at all, so no bibliography entry of the article is transcribed and no reference is redistributed. Reference adaptation: none was needed and none was made; every reference inside the delivered unit resolves inside it, so no unresolved reference or citation is produced. Printed slips kept literal: no printed word, symbol, hypothesis or number of the counted statement or of its proof was altered. Layout and class: the article's own class is \texttt{article} with packages including graphicx, inputenc, geometry, authblk, cite, algorithm, algpseudocode, dsfont, amsmath, amssymb, amsthm, xcolor, hyperref, cleveref; the wrapper is the lane's pinned \texttt{amsart} editorial wrapper at 10pt on A4, which restates the theorem-like environments the retained text needs and adds the article's own macro definitions that the retained text needs (\texttt{\textbackslash CC}, \texttt{\textbackslash RR}, \texttt{\textbackslash Trop}, \texttt{\textbackslash Tropw}, \texttt{\textbackslash initial}, \texttt{\textbackslash inw}), with extra wrapper packages (bm, bbm, dsfont, xspace, cancel, mathtools, cleveref). The delivered numbering is the unit's own. Assets: no retained span contains a figure environment, a graphics-inclusion command, a PDF-page inclusion, a table or a TikZ picture; no image file is copied into this package, the package declares no asset dependency and no image file is redistributed with it. Licence: version arXiv:2408.15719v2 of this article is distributed under the Creative Commons Attribution 4.0 International licence, as recorded in the arXiv licence metadata for this exact version and shown on the abstract page https://arxiv.org/abs/2408.15719v2. A full rights notice is delivered at licenses/arxiv-2408.15719-CC-BY-4.0.txt. Characters: every retained excerpt is pure ASCII, so the delivered engine, XeLaTeX with its default Latin Modern text font, reports no missing character.
\begin{proposition}
\label{Prop::CharWeakSigned}
Let  $V \subseteq (\CC^*)^n$ be a very affine algebraic variety defined by an ideal $I \subseteq \mathcal{C}[x_1^\pm, \dots , x_n^\pm]$. For $w \in \mathbb{R}^n$ the following are equivalent
\begin{itemize}
    \item[(i)] $ w \in \Trop_\CC^+(V)$,
    \item[(ii)] $\exists u \in \mathbb{R}^n \colon V\big( \initial_u(\initial_w(I)) \big) \cap \widetilde{\mathcal{R}}^{n}_{>0} \neq \emptyset$,
    \item[(iii)] $\initial_w(I) \cap \widetilde{\mathcal{R}}_{\geq 0}[x_1, \dots , x_n] = \{ 0 \} $.
\end{itemize}
 Furthermore, the set of points $w \in \mathbb{R}^n$ with $V\big( \initial_w(I) \big) \cap \widetilde{\mathcal{R}}^{n}_{>0} \neq \emptyset $ forms a dense subset of $\Trop_\CC^+(V)$.
\end{proposition}

\begin{proposition}
\label{prop: positive part of toric vars}
For $Y_A =  \varphi_A \left( (\CC^*)^n \right)  \subseteq (\mathcal{C}^*)^r$ a very affine toric variety defined by a matrix $A \in \mathbb{Z}^{n \times r}$ of rank $n \leq r$, we have
\[ Y_A \cap \RR^r_{>0} = \varphi_A(\RR^n_{>0}).\]
\end{proposition}

\begin{lemma}
\label{lemma: induced map of initial degs}
For every weight vector $w \in \mathbb{R}^n$,
the monomial map $\varphi_A$ induces a monomial map on the initial degenerations
\[
\overline{\varphi_A}\colon \inw (X) \longrightarrow \initial_{A^{\top}w} (Y).
\]
Furthermore, if $X = \varphi_A^{-1}(Y)$, then
the schemes $\inw X $ and 
$\overline{\varphi_A}^{-1}(\initial_{A^\top w} (Y))$ agree.
\end{lemma}

\begin{proposition}
\label{prop: functoriality of monomial maps}
let $I \subseteq \CC[x_1^\pm, \dots, x_n^\pm]$ be a prime ideal defining a variety $X \subseteq \mathbb{G}^n_\CC$ and let
$J \subseteq \CC[x_1^\pm, \dots, x_r^\pm]$ be a prime ideal defining $Y \subseteq \mathbb{G}^r_\CC$
with $\varphi_A(X) \subseteq Y$.

    The linear images of the positive and of the weakly positive tropicalization of $X$ are contained in the respective tropicalizations of $Y$:
    \begin{align*}
            A^\top \Trop^+(X) \subseteq \Trop^+(Y), \qquad A^\top\Tropw(X) \subseteq \Tropw(Y).
    \end{align*}
    If $\varphi_A(X \cap \RR^n_{> 0 }) = Y\cap\RR^r_{> 0}$,
    then the first inclusion is an equality.
    Furthermore, if both $Y$ is contained in the image of $\varphi_A$,
    and $X = \varphi_A^{-1}(Y)$,
    then the second inclusion is also an equality.
\end{proposition}

\begin{proof}
    First we show the inclusion $A^\top \Trop^+(X) \subseteq \Trop^+(Y)$. The proof for the inclusion $A^\top \Tropw(X) \subseteq \Tropw(Y)$ is analogous.
    Let $x$ be any point in $X \cap \mathcal{R}^n_{>0}$ and let $w = \nu(x) \in \Trop^+(X)$. 
    Direct computation shows $\nu(\varphi_A(x)) = A^\top w$, giving the inclusion.

    We now assume $\varphi_A(X \cap \RR^n_{> 0 }) = Y\cap\RR^r_{> 0}$ and show the inclusion  $\Trop^+(Y) \subseteq A^\top\Trop^+(X)$.
    Let $y$ be in $Y \cap \RR^r_{> 0}$. By assumption there exists a point $x \in X \cap \RR^n_{>0}$ with $\varphi_A(x) = y$. Then $A^\top \nu(x) = \nu(y)$, finishing this part of the proof.

    We now assume $\varphi_A(X \cap \RR^n_{> 0 }) = Y\cap\RR^r_{> 0}$ 
    and $X = \varphi_A^{-1}(Y)$
    in order to finally show the desired equality $A^\top\Tropw(X) = \Tropw(Y)$ for the weakly positive tropicalization.
    Let $y \in Y, \ u  = \nu(y)$ such that
    $\initial_{u} (Y) $ contains a real point with positive entries.
    By \Cref{Prop::CharWeakSigned} these points form a dense subset of $\Tropw(\varphi_A(Y))$.
     Let $x \in \varphi_A^{-1}(y)$ and let $w = \nu(x)$.
    By \Cref{lemma: induced map of initial degs} the initial degeneration $\inw(X)$ is the preimage of $\initial_{u} (Y)$ under a monomial map. By \Cref{prop: positive part of toric vars} $\inw(X)$ contains a positive point, showing $w \in \Tropw(X)$.
    This shows the inclusion $\Tropw(Y) \subseteq A^\top\Tropw(X)$ and finishes the proof.
\end{proof}

\end{document}
